
\subsubsection{2.1 Self-Repair and Iterative Refinement}

LLM-based self-repair has emerged as a paradigm for improving code generation accuracy. Prior work systematically evaluated self-repair on HumanEval, finding improvements of 4.9--17.1 percentage points over single-shot generation. Repair success varies by error type, with assertion errors among the more challenging categories. Multi-agent collaborative approaches have achieved 77.13\% Pass@1, representing current state-of-the-art performance. However, these works focus on whether repair works rather than why certain feedback enables better repair.

Debugging decay---substantial capability loss after multiple repair attempts---has been identified, suggesting that feedback quality in early attempts is crucial. This motivates a focus on single-attempt success and understanding what makes initial feedback effective.

\subsubsection{2.2 Trace-Based Debugging}

Prior work has demonstrated that runtime traces outperform error messages for code repair, using debugging-style execution to provide richer context. Self-Debug introduced execution trace methodology for self-repair. While these works establish the superiority of traces over messages, they do not decompose which aspects of traces drive improvement. The AS framework attempts this decomposition.

\subsubsection{2.3 Type-Guided Generation}

Type system internalization has been shown to improve functional correctness by providing structured constraints. This suggests that specific, actionable information (types) helps LLMs generate correct code. The AS framework generalizes this insight: type information contributes to AS\_loc (type error locations) and AS\_state (type constraints as state).

\subsubsection{2.4 Feedback Informativeness Gap}

Prior work treats feedback signals as categorical (trace vs. message vs. static analysis) rather than decomposing their information content. Combined compiler and symbolic execution feedback has been used, but without measuring which components drive improvement. Agentic approaches have achieved 42.3\% solve rate with 11.8 average iterations, demonstrating that iteration quantity alone is insufficient.

The present work fills this gap by proposing measurable AS components and testing their extractability. The negative result---that regex extraction fails on assertion-dominated benchmarks---reveals a boundary condition that prior work implicitly assumed away.

